% !TeX program = lualatex
% Timelines for the Arlington County Board Structure and Performance report.
%
% An intermediate product, not part of the report. The dated rows the report's
% prose is written from live here, one block per timeline, grouped by the
% section of the report they belong to, so a timeline can leave the paper
% without being lost. Compiled the same way as the paper (lualatex, biber) and
% drawing on the same paper/sources.bib.
\documentclass[11pt]{article}

\usepackage{booktabs}
\usepackage[margin=1in]{geometry}
\usepackage{hyperref}

\usepackage{titlesec}
\titleformat{\section}{\normalfont\normalsize\bfseries}{}{0em}{}
\titlespacing*{\section}{0pt}{16pt}{6pt}

\usepackage{fontspec}
\setmainfont{Lato}[
  Path           = ../style/fonts/,
  Extension      = .ttf,
  UprightFont    = *-Regular,
  BoldFont       = *-Bold,
  ItalicFont     = *-Italic,
  BoldItalicFont = *-BoldItalic,
]

\usepackage[notes, backend=biber]{biblatex-chicago}
\addbibresource{sources.bib}
\AtEveryCitekey{\clearfield{url}\clearfield{urldate}\clearfield{urlyear}\clearfield{urlmonth}\clearfield{urlday}}

% A footnote inside a tabular is silently dropped by LaTeX: the mark prints
% and the note never appears. Timeline rows carry citations, so make tabular
% save and emit them.
\usepackage{footnote}
\makesavenoteenv{tabular}

% A timeline is a dated list: a narrow date column and a wrapping one beside it.
\newcommand{\timeline}[1]{%
  \begin{center}
  \begin{tabular}{@{}p{0.17\textwidth}p{0.77\textwidth}@{}}
    \toprule
    #1
    \bottomrule
  \end{tabular}
  \end{center}}

% Every section starts on a new page, as in the paper.
\let\oldsection\section
\renewcommand{\section}{\clearpage\oldsection}

% A comma between two footnote marks that stand together.
\newcommand{\fnsep}{\textsuperscript{,}}

\title{Timelines for the Arlington County Board report}
\author{Sally Hudson and Alex Keena}
\date{\today}

\begin{document}

\maketitle

\section{Board Structure}

\timeline{
  1846--1847 & Congress returns Alexandria County to Virginia, with the assent of
  its voters in a referendum.\autocite{usstat1846retrocession} \\[2pt]
  1847--1870 & A county court of elected justices of the peace governs the county,
  and its work is mainly judicial. No elected body does what the Board now
  does.\autocite[art.~VI, secs.~25--27]{vaconstitution1851} \\[2pt]
  1870--1931 & Virginia's constitution creates boards of supervisors for counties
  statewide. Three magisterial districts --- Arlington, Jefferson and
  Washington --- elect one supervisor each.\autocite[art.~VII,
  sec.~2]{vaconstitution1869} \\[2pt]
  1930 & The county adopts the County Manager Plan by
  referendum.\autocite{samuel2026} \\[2pt]
  1931--1932 & The first at-large Board of five is elected in November 1931 and
  seated in January 1932, in place of three district seats. \\[2pt]
  1938 & Voters approve staggered terms, 1,539 to 1,487 in the county's tally;
  from the Board elected in 1939, one member is chosen in most years and two in
  the fourth.\autocite[11]{arlingtonelections2021} \\
}

\section{Board Size}

\timeline{
  1930--1950 & Arlington's population grows more than fivefold.
  % The deck gives a cause (the federal workforce, the New Deal and World War II)
  % that no source held supports yet: population-growth-causes
  \\[2pt]
  1952--1958 & The General Assembly lets Arlington voters decide, on a petition
  of 200 voters, whether to raise the Board from five seats to seven and
  to elect members every two years. No petition is filed. It reopens the
  window for the election question in 1954 and again in 1958, without the
  seven-seat question, and again no petition is
  filed.\autocite{vaacts1952c591,vaacts1954c151,vaacts1958c207} \\
}

\section{Gender}

\timeline{
  1932 & Elizabeth Magruder takes a seat, the first woman on the Board, in the
  year the Board expands to five and moves to at-large election. \\[2pt]
  1938 & Magruder is the first woman to chair the Board, and chairs three years
  running.\autocite{arlingtonva2026members} \\[2pt]
  1938 & The Arlington County Woman's Democratic Club opposes staggered terms,
  holding that they would make it ``virtually impossible for a woman to be
  elected to the board.''\autocite{sun1938referendum} \\[2pt]
  1959--1973 & 15 consecutive years with no woman on the Board. \\[2pt]
  1974 & Ellen Bozman takes a seat, the first woman on the Board in 15
  years. \\[2pt]
  1983, 2024 & The only two years in which women hold a majority of the seats:
  Bozman, Dorothy Grotos and Mary Margaret Whipple in 1983; Libby Garvey, Maureen
  Coffey and Susan Cunningham in 2024. \\
}

\section{Race and Ethnicity}

\timeline{
  1870--1890 & A majority-Black county until the 1890 census, which counts it
  at 49.9 percent. The Board is created in 1870. Freedman's Village, the
  settlement of formerly enslaved people on the Lee estate, lies inside
  Arlington district and counts 338 residents in
  1890.\autocite{census1890} \\[2pt]
  1890--1950 & The White population grows from 2,135 to 128,780 and the Black
  share falls from 49.9 percent to 4.8; the Black population grows from 2,123 to
  6,517. % waits on: population-growth-causes
  \\[2pt]
  1980--present & Hispanic residents grow from 8,863 to 37,362 and Asian and
  Pacific Islander residents from 6,792 to 27,353, from the first census to ask
  every household its Hispanic origin. \\
}

\timeline{
  \multicolumn{2}{@{}l}{\textit{1871--1892 --- Black representation under the
  district system}} \\[4pt]
  1888 & The county court issues a rule against Tibbett Allen, the Board's only
  Black member, to show cause why he should not be removed for
  non-residence; he resigns, and the court never
  rules.\autocite{alexandriagazette1888rule} \\[2pt]
  1889--92 & Walter G. Willson holds the Washington district seat, the last
  Black member until 1988.\autocite{alexandriagazette18890525p3,alexandriagazette18910529p3,gazette1892wilsondeath} \\[4pt]
  \multicolumn{2}{@{}l}{\textit{1893--1919 --- Suppressing Black
  representation}} \\[2pt]
  1894 & The Walton Act makes the ballot official and bars any other as a legal
  vote. A voter must strike through the name of every candidate he does not
  want, with a line reaching three-fourths of the name's length or the stroke
  does not count, and an office left with more than one name standing is void.
  Only a constable the electoral board appoints may read or mark a ballot for a
  voter who cannot do it himself, and giving out a copy of the ballot is a
  misdemeanor.\autocite{waltonact1894}
  \\[2pt]
  1901--02 & The Virginia Constitutional Convention writes a poll tax and a
  literacy test into the constitution. \\[4pt]
  \multicolumn{2}{@{}l}{\textit{1920--1940 --- Replacing district
  representation with countywide majority contests}} \\[2pt]
  1922 & \textit{Bennett v. Garrett}: the Supreme Court of Appeals affirms the
  dismissal of Clarendon's petition to incorporate as a town, holding that where
  proposed limits are carved out of a thickly settled section the ``general good
  of the community'' the statute protects ``must be construed as embracing the
  whole section.''\autocite{bennettvgarrett1922} \\[2pt]
  1930--31 & A referendum adopts at-large election in place of magisterial
  districts, and the first at-large election follows in November 1931,
  changing who the constituency is. \\[2pt]
  1938 & Staggered terms, approved 1,539 to 1,487 in the county's tally, so
  most elections fill a single seat --- changing the threshold, not the
  constituency.\autocite[11]{arlingtonelections2021} \\
}

\timeline{
  \multicolumn{2}{@{}l}{\textit{1939--1987 --- No Black member serves}} \\[2pt]
  1974--75 & \textit{Vollin v. Kimbel}: eight Arlington residents from Black,
  Hispanic and lower-income neighborhoods sue to have at-large election
  declared unconstitutional; the district court dismisses the suit in 1974 and
  the Fourth Circuit affirms in 1975.\autocite{vollin1974} \\[2pt]
  1976 & More than 200 voters petition to put district election to a
  referendum; the Virginia Supreme Court denies the petition because the County
  had already adopted the County Manager Plan.\autocite{vollinelectoralboard}
  \\[4pt]
  \multicolumn{2}{@{}l}{\textit{1988--present --- Members of color return}}
  \\[2pt]
  1988-- & Newman, Monroe, Tejada, Dorsey, Talento, Spain. \\
}

\section{Project Scope: Local legal authority}

\timeline{
  1882 & \textit{Kirkham v. Russell}: refusing mandamus to seven men whom
  Petersburg's outgoing common council had elected to city offices three days
  before 11 of its 24 members left office, the Supreme Court of
  Appeals cites Dillon on Municipal Corporations for the rule that a charter
  grants only the powers its words comprehend or necessarily
  imply.\autocite{kirkhamvrussell1882} \\[2pt]
  1896 & \textit{City of Winchester v. Redmond}: refusing to pay a reward
  Winchester's council had promised for catching incendiaries, the court
  quotes Dillon directly for the rule that now carries his name. Schragger and
  Retzloff name Kirkham the case in which Virginia adopted the rule in full and
  Winchester the one that quotes it
  outright.\autocite{winchestervredmond1896}\fnsep\autocite[195--96]{schragger2022}
  \\[2pt]
  1968--71 & Virginia's constitutional revision commission proposes letting a
  charter county or city act on anything state law does not deny it; the
  General Assembly strikes the provision before the new constitution reaches
  voters. The opposition is the Virginia Municipal League's --- ``somewhat
  counterintuitive'', Schragger and Retzloff call it, since the Virginia
  Association of Counties had proposed the opposite in 1968 and reverses its
  later hostility only in its 2020 legislative
  program.\autocite[200]{schragger2022} \\[2pt]
  1971 & The General Assembly conforms the County Manager Plan's statutory
  text to the new Constitution: the sections setting the board's membership,
  the manager's powers and the board's authority to abolish offices each
  replace a cross-reference to the old constitution's section 110 with Article
  VII, Section 4, and nothing else changes.\autocite{vaacts1971exsess} \\[2pt]
  1977 & \textit{Commonwealth v. Arlington County Board}: the Virginia Supreme
  Court voids the County's collective bargaining agreements with its own
  employees and, on a question of first impression for every local governing
  body, bars public-sector bargaining statewide.\autocite{commonwealthvarlington1977}
  \\[2pt]
  1985 & \textit{County Board of Arlington v. Brown}: the Board accepts a letter
  of intent for a 75-year lease of 6.4 acres of the courthouse
  parking lot to a private developer and directs its County Manager to sign;
  he refuses, the Board sues to compel him, and the Virginia Supreme Court
  denies the writ because a county has no power to
  lease.\autocite{brownvarlington1985} \\[2pt]
  2000 & \textit{Arlington County v. White}: the County extends its health plan
  to employees' domestic partners, and the Virginia Supreme Court holds the
  extension beyond the County's authority.\autocite{arlingtonvwhite2000} \\
}

\section{The Board and the Manager}

\timeline{
  1930 & The County Manager Plan has the Board appoint a manager, who holds the
  county's administrative and executive powers, including appointments; the
  Board may not move a budget allocation by more than 10
  percent.\autocite{vaacts1930c167} \\[2pt]
  1937 & The Board delegates the appointment of department heads to the
  Manager.\autocite{sun1951rulingboard} \\[2pt]
  1951 & The Board takes the power back by resolution in January; in March a
  circuit judge rules that it had no right to give it up.%
  \autocite{dailysun1951spicer,sun1951rulingboard} \\[2pt]
  1952 & The General Assembly bars Board members from directing or requesting
  appointments; in November voters approve an indefinite term for the Manager
  and appointment of department heads by the Manager.%
  \autocite{vaacts1952c443,vaacts1952c198} \\[2pt]
  1982 & The Board's personnel-interference bar loses its criminal penalty and the
  Board may discuss appointments with the Manager.
  % Unconfirmed: whether 1982 c. 108 or the 1997 recodification made the change
  % (session-acts-1982-c108).
  \\
}

\clearpage
\printbibliography[title={Works cited}]

\end{document}
